% !TeX root = ../main.tex
\chapter{Introducción}
\section{Fundamentación}

La física acústica es fundamental en la formación de un fonoaudiólogo, ya que permite comprender los principios físicos del sonido, su propagación y las características de las ondas sonoras, aspectos esenciales para evaluar y tratar trastornos de la comunicación. Con este conocimiento, el fonoaudiólogo puede entender cómo se generan, transmiten y perciben los sonidos, lo cual es crucial en la rehabilitación auditiva y del habla. Además, la física acústica ofrece herramientas para analizar la calidad del sonido y su impacto en el sistema auditivo, lo que permite mejorar las intervenciones terapéuticas y el uso de tecnologías como los audífonos.

\section{Objetivos generales y específicos}

\subsection{Objetivos generales}

\begin{itemize}
     \setlength\itemsep{0.3em}        
	\item Brindar al estudiante los elementos necesarios para alcanzar una perspectiva global del tema.
	\item Comprender y manejar los conceptos básicos de acústica y psicoacústica.

\end{itemize}   

\subsection{Objetivos específicos}  

\begin{itemize}
	\item Conocer los alcances de la acústica y su relación con la Fonoaudiología.
	\item Identificar las incumbencias del profesional en esta área.
	\item Interpretar la naturaleza del sonido y su carácter objetivo-subjetivo.
	\item Interpretar y operar con las magnitudes básicas empleadas para cuantificar los fenómenos sonoros.
	\item Entender los procesos físicos que tienen lugar en el oído humano.
        \item Caracterizar las distintas enfermedades auditivas y conocer los métodos de diagnostico de cada una.
        \item Comprender los fundamentos de los principales fenómenos psicoacusticos.
        \item Asimilar el concepto de ruido, conocer su modo de medición y sus efectos adversos.
        \item Aprender a interpretar gráficos de respuesta en frecuencia.

\end{itemize}

%%%%%%%%%%%%%%%%%%%%%%%%%%%%%%%%%%%%%%%%%%%%%%%%%%%%%%%%%%%%%%%%%%%%%%%%%%%%%%%%%%%%%%%%%%%%%%%%%%
